% ============================================================
%  AGRITWIN — Rapport de Projet — Prix de l'Innovation 2026
%  Polytech Innovation Day — École Polytechnique de Thiès
%  Mise en page calquée sur le modèle fourni (Yoonwi).
% ============================================================

\documentclass[12pt, a4paper]{article}

% ── Encodage & langue ──────────────────────────────────────
\usepackage[utf8]{inputenc}
\usepackage[T1]{fontenc}
\usepackage[french]{babel}

% ── Typographie ────────────────────────────────────────────
\usepackage[stretch=10,shrink=10,expansion=false]{microtype}
\usepackage{setspace}
\setstretch{1.3}

% ── Mise en page ───────────────────────────────────────────
\usepackage[top=2.5cm, bottom=2.8cm, left=3cm, right=3cm]{geometry}

% ── Couleurs ───────────────────────────────────────────────
\usepackage{xcolor}
\definecolor{primary}{HTML}{14532D}
\definecolor{accent}{HTML}{15803D}
\definecolor{lightbg}{HTML}{F0FDF4}
\definecolor{rule}{HTML}{CBD5E1}
\definecolor{dark}{HTML}{0F172A}
\definecolor{muted}{HTML}{5B6A7D}
\definecolor{mygreen}{HTML}{15803D}
\definecolor{techblue}{HTML}{0284C7}

% ── Graphiques ─────────────────────────────────────────────
\usepackage{graphicx}
\usepackage{relsize}
\usepackage{tikz}
\usetikzlibrary{positioning, arrows.meta, shapes.geometric}

% ── Tableaux ───────────────────────────────────────────────
\usepackage{booktabs}
\usepackage{tabularx}
\usepackage{colortbl}
\usepackage{array}

% ── Boîtes légères ─────────────────────────────────────────
\usepackage{mdframed}
\mdfdefinestyle{primary}{
  linecolor=primary, linewidth=1.2pt,
  backgroundcolor=primary!5,
  leftline=true, rightline=false, topline=false, bottomline=false,
  innerleftmargin=12pt, innertopmargin=8pt, innerbottommargin=8pt,
  skipabove=8pt, skipbelow=8pt
}
\mdfdefinestyle{accent}{
  linecolor=accent, linewidth=1.2pt,
  backgroundcolor=accent!7,
  leftline=true, rightline=false, topline=false, bottomline=false,
  innerleftmargin=12pt, innertopmargin=8pt, innerbottommargin=8pt,
  skipabove=8pt, skipbelow=8pt
}
\mdfdefinestyle{neutral}{
  linecolor=rule, linewidth=0.8pt,
  backgroundcolor=lightbg,
  leftline=true, rightline=false, topline=false, bottomline=false,
  innerleftmargin=12pt, innertopmargin=8pt, innerbottommargin=8pt,
  skipabove=8pt, skipbelow=8pt
}

% ── En-têtes & pieds de page ───────────────────────────────
\usepackage{fancyhdr}
\pagestyle{fancy}
\fancyhf{}
\fancyhead[L]{\small\textcolor{primary}{\textbf{AGRITWIN}}%
  \quad\textcolor{rule}{|}\quad\small\textcolor{muted}{Prix de l'Innovation 2026}}
\fancyhead[R]{\small\textcolor{muted}{Polytech Innovation Day — EPT}}
\fancyfoot[C]{\small\textcolor{muted}{\thepage}}
\renewcommand{\headrulewidth}{0.4pt}
\renewcommand{\headrule}{\hbox to\headwidth{\color{primary}%
  \leaders\hrule height \headrulewidth\hfill}}

% ── Sections ───────────────────────────────────────────────
\usepackage{titlesec}
\titleformat{\section}
  {\color{primary}\large\bfseries}
  {\thesection.}{0.6em}{\MakeUppercase}
  [\vspace{2pt}\textcolor{accent}{\rule{\linewidth}{1.5pt}}]

\titleformat{\subsection}
  {\color{dark}\normalsize\bfseries}
  {\thesubsection.}{0.5em}{}

\titleformat{\subsubsection}
  {\color{primary}\normalsize\itshape\bfseries}
  {}{0em}{}

\titlespacing{\section}{0pt}{20pt}{8pt}
\titlespacing{\subsection}{0pt}{12pt}{4pt}
\titlespacing{\subsubsection}{0pt}{8pt}{2pt}

% ── Liens ──────────────────────────────────────────────────
\usepackage[
  colorlinks=true, linkcolor=primary,
  urlcolor=accent, citecolor=mygreen,
  pdftitle={AgriTwin — Prix de l'Innovation 2026},
  pdfauthor={Equipe AgriTwin — EPT}
]{hyperref}

% ── Listes ─────────────────────────────────────────────────
\usepackage{enumitem}
\setlist[itemize]{
  leftmargin=1.5em,
  label=\textcolor{accent}{\small$\bullet$},
  itemsep=3pt, topsep=4pt
}
\setlist[enumerate]{itemsep=3pt, topsep=4pt}

% ── Table des matières ─────────────────────────────────────
\usepackage{tocloft}
\renewcommand{\cftsecfont}{\bfseries\color{primary}}
\renewcommand{\cftsecpagefont}{\bfseries\color{primary}}
\renewcommand{\cftsubsecfont}{\color{dark}}
\renewcommand{\cftsubsecpagefont}{\color{dark}}

% ══════════════════════════════════════════════════════════
\begin{document}
% ══════════════════════════════════════════════════════════

% ──────────────────────────────────────────────────────────
%  PAGE DE GARDE
% ──────────────────────────────────────────────────────────
\thispagestyle{empty}

\vspace*{1cm}

\begin{center}

  % ── Emplacement du logo : remplacer le cadre par ────────
  %    \includegraphics[width=3cm]{agritwin-logo.png}
  \begin{tikzpicture}
    \draw[rule, line width=0.6pt, rounded corners=4pt] (0,0) rectangle (3,3);
    \node[text=muted, font=\small\itshape] at (1.5,1.5) {Logo};
  \end{tikzpicture}
  % ───────────────────────────────────────────────────────

  \vspace{0.2cm}
  {\color{rule}\rule{5cm}{0.4pt}}

  \vspace{1.8cm}
  {\Huge\textbf{\textcolor{primary}{\scalebox{2.5}{AGRITWIN}}}}

  \vspace{0.5cm}
  {\large\textcolor{accent}{Le jumeau numérique qui transforme une parcelle en preuve de crédit}}

  \vspace{0.4cm}
  \textcolor{rule}{\rule{10cm}{0.4pt}}

  \vspace{0.6cm}
  {\normalsize\textcolor{muted}{\itshape
    \og{}L'ingénierie face aux enjeux du financement de l'agriculture familiale\fg{}}}

  \vspace{0.3cm}
  {\small\textbf{\textcolor{dark}{Prix de l'Innovation 2026 — Polytech Innovation Day}}}

  \vspace{2.6cm}

  \begin{tabular}{ll}
    \textcolor{primary}{\textbf{[Prénom Nom]}} & [Formation] \\[5pt]
    \textcolor{primary}{\textbf{[Prénom Nom]}} & [Formation] \\[5pt]
    \textcolor{primary}{\textbf{[Prénom Nom]}} & [Formation] \\[5pt]
    \textcolor{primary}{\textbf{[Prénom Nom]}} & [Formation] \\
  \end{tabular}

  \vspace{1.3cm}
  \textcolor{rule}{\rule{10cm}{0.4pt}}

  \vspace{0.4cm}
  {\small\textcolor{muted}{École Polytechnique de Thiès \quad — \quad Juin 2026}}

\end{center}

\newpage

% ──────────────────────────────────────────────────────────
%  TABLE DES MATIÈRES
% ──────────────────────────────────────────────────────────
\tableofcontents
\newpage

% ══════════════════════════════════════════════════════════
\section{Fiche d'identification du projet}
% ══════════════════════════════════════════════════════════

\begin{center}
\begin{tabularx}{\linewidth}{>{\bfseries\color{primary}}p{4.2cm} X}
  \toprule
  \rowcolor{lightbg}
  \multicolumn{2}{l}{\textbf{\textcolor{primary}{Dossier de candidature — Prix de l'Innovation 2026}}} \\
  \midrule
  Titre du projet & \textbf{AgriTwin} — Jumeau numérique de parcelle et Score Vert opposable \\
  \addlinespace
  Porteurs du projet & [Prénom Nom] · [Prénom Nom] · [Prénom Nom] · [Prénom Nom] \\
  \addlinespace
  Établissement & École Polytechnique de Thiès (EPT) \\
  \addlinespace
  Formation d'appartenance & [GEM / GIT / GA…] \\
  \addlinespace
  Encadrant académique & [Nom de l'encadrant, le cas échéant] \\
  \addlinespace
  Contact & [adresse électronique] · [téléphone] \\
  \addlinespace
  Domaine & Agriculture numérique, finance rurale, traçabilité \\
  \addlinespace
  Niveau de maturité & Prototype fonctionnel (TRL 4) — données simulées \\
  \addlinespace
  Prototype & [URL du prototype à compléter] \\
  \bottomrule
\end{tabularx}
\end{center}

\begin{mdframed}[style=neutral]
\textbf{À compléter avant dépôt.} Les champs entre crochets sont volontairement laissés vides : remplacer les noms, la formation, l'encadrant, le contact et l'URL. Le cadre du logo en page de garde s'ouvre avec \texttt{\textbackslash includegraphics}.
\end{mdframed}

% ══════════════════════════════════════════════════════════
\section{Résumé exécutif}
% ══════════════════════════════════════════════════════════

\begin{mdframed}[style=accent]
\textbf{AgriTwin} transforme les relevés d'un boîtier scellé, les observations
satellite et les pièces d'un dossier en un \textbf{Score Vert} vérifiable, qu'un
exploitant agricole, une institution de microfinance et un agent de conformité
consultent dans la même application — sans base de données distante, sans API
payante et sans GPS embarqué sur le matériel.
\end{mdframed}

\vspace{0.4cm}

Le problème n'est pas l'absence de données agricoles, mais l'\textbf{absence de preuve
opposable}. Un producteur de Bambey peut déclarer 1,5 hectare d'oignon ; la banque, elle,
n'a aucun moyen de vérifier ni la surface, ni la culture, ni la réalité des relevés qui
justifieraient un crédit d'intrants. La vérification coûte plus cher que le crédit
lui-même. Le résultat est connu : les petits producteurs restent hors du système financier
formel, ou empruntent à des taux qui absorbent leur marge.

AgriTwin renverse la logique : la preuve est produite \emph{au moment de la mesure},
signée, chaînée, puis restituée sous forme de score décomposé et d'un dossier exportable.
Le prototype présenté ici couvre l'intégralité du parcours — déclaration géolocalisée,
contrôle de conformité, surveillance continue, détection de fraude, audit cryptographique,
décision bancaire — et fonctionne entièrement dans un navigateur.

\vspace{0.5cm}

\begin{center}
\begin{tabular}{p{4.2cm} p{4.2cm} p{4.2cm}}
  \toprule
  \rowcolor{lightbg}
  \textbf{\textcolor{primary}{Problème}} &
  \textbf{\textcolor{primary}{Solution}} &
  \textbf{\textcolor{primary}{Impact visé}} \\
  \midrule
  Déclarations invérifiables, relevés non opposables, coût de contrôle prohibitif. &
  Boîtier scellé, chaîne SHA-256, Score Vert décomposé, portail prêteur sous consentement. &
  Dossiers instruisibles à distance, délais de décision réduits, accès au crédit élargi. \\
  \bottomrule
\end{tabular}
\end{center}

% ══════════════════════════════════════════════════════════
\section{Présentation synthétique de la problématique}
% ══════════════════════════════════════════════════════════

\subsection{Un crédit agricole bloqué par l'incertitude}

Le financement de campagne repose sur trois vérifications que personne ne peut faire à
coût raisonnable : \textbf{la parcelle existe-t-elle} avec la surface déclarée,
\textbf{la culture est-elle réellement conduite} comme annoncé, et \textbf{les résultats
techniques} (eau, azote, santé du sol) sont-ils ceux qui justifient le montant demandé.

\begin{mdframed}[style=primary]
Au Sénégal, l'obstacle n'est pas le manque de moyens de production, mais \textbf{le manque
de preuve vérifiable} sur ces moyens. Un agent qui doit se déplacer, mesurer et attester
coûte plus cher que la marge du dossier : la vérification n'a pas lieu, et le crédit est
refusé ou renchéri.
\end{mdframed}

\subsection{Quatre frictions observées sur le terrain}

\begin{tabularx}{\linewidth}{>{\bfseries\color{primary}}l X}
  \toprule
  \rowcolor{lightbg}
  Situation & Friction vécue \\
  \midrule
  Déclaration de surface
    & Surface annoncée non concordante, contour absent, pas de photo horodatée \\
  \addlinespace
  Contrôle des pièces
    & Pièces papier, copies non vérifiables, dossiers perdus entre deux guichets \\
  \addlinespace
  Suivi de campagne
    & Aucun historique de l'irrigation ni des apports : la banque ne peut pas expliquer un échec \\
  \addlinespace
  Aléa et fraude
    & Un relevé peut être fabriqué ; rien ne permet de distinguer une panne d'une falsification \\
  \bottomrule
\end{tabularx}

\subsection{Pourquoi maintenant}

\begin{itemize}
  \item Le matériel de mesure agricole est devenu accessible : une sonde 4-en-1
        (humidité, température, pH, conductivité) coûte environ 54\,000 FCFA rendue à Dakar.
  \item Le satellite est gratuit et régulier : Sentinel-2 fournit le NDVI, Sentinel-1
        l'humidité de surface, tous les cinq à six jours.
  \item Le paiement mobile (Wave, Orange Money) rend le décaissement et le remboursement
        traçables sans infrastructure bancaire locale.
  \item Les institutions de microfinance et les coopératives cherchent des dossiers
        \emph{instruits} : la contrainte est la qualité du dossier, pas le volume de demandes.
\end{itemize}

% ══════════════════════════════════════════════════════════
\newpage
\section{Description de la solution innovante}
% ══════════════════════════════════════════════════════════

\subsection{Principe : la preuve naît avec la mesure}

AgriTwin est un \textbf{jumeau numérique de parcelle}. Chaque parcelle déclarée possède un
contour géographique, une culture, un boîtier scellé et un historique de relevés. Chaque
relevé est sérialisé de façon canonique, haché en SHA-256 avec l'empreinte du relevé
précédent, puis chaîné. Toute réécriture casse la chaîne et devient visible.

Cette chaîne alimente un \textbf{Score Vert sur 100}, décomposé en quatre composantes que
la banque peut discuter une par une. Le score n'est donc pas une note opaque : c'est un
argumentaire chiffré, avec ses poids, ses seuils et ses limites publiés dans l'application.

\subsection{Le parcours complet, en cinq temps}

\begin{center}
\begin{tikzpicture}[
  step/.style={
    rectangle, rounded corners=5pt, minimum width=11.5cm, minimum height=1.05cm,
    draw=primary, fill=primary!6, thick,
    text centered, font=\small\bfseries\color{primary}
  },
  arr/.style={-{Stealth[length=5pt]}, thick, accent}
]
  \node[step] (a) {1. Déclarer : contour, culture, pièces justificatives};
  \node[step, below=0.45cm of a] (b) {2. Contrôler : conformité des pièces et du contour};
  \node[step, below=0.45cm of b] (c) {3. Mesurer : boîtier scellé + Sentinel-1/2, 90 jours};
  \node[step, below=0.45cm of c] (d) {4. Prouver : chaîne SHA-256, Score Vert décomposé};
  \node[step, below=0.45cm of d] (e) {5. Décider : portail prêteur, PDF et archive vérifiable};

  \foreach \x/\y in {a/b, b/c, c/d, d/e}
    \draw[arr] (\x) -- (\y);
\end{tikzpicture}
\end{center}

\subsection{Quatre piliers fonctionnels}

\subsubsection{Déclaration et pièces vérifiées}

Parcours en quatre étapes : identité de l'exploitation, localisation et tracé du contour
sur fond satellite (surface recalculée géométriquement), pièces justificatives, puis
récapitulatif. Les fichiers sont contrôlés sur leur \textbf{signature binaire réelle} — pas
sur leur extension ni sur le type MIME déclaré — avec des plafonds de 10\,Mo par pièce,
25\,Mo par dossier et 4 photos. Un dossier soumis ou certifié repasse automatiquement en
brouillon dès qu'il est modifié.

\subsubsection{Surveillance continue et alertes bilingues}

Le boîtier fournit humidité à 15 et 35\,cm, températures sol et air, pluie, conductivité
et pH, état du scellé, accéléromètre, batterie et niveau radio. Huit règles d'alerte
couvrent l'incohérence sonde/satellite, la rupture de scellé, la chaleur et le stress
hydrique, la salinité et les signalements phytosanitaires. Les messages existent en
\textbf{français et en wolof} : le texte wolof est affiché avec sa traduction, et un
repère sonore accompagne la lecture.

\subsubsection{Score Vert et preuve cryptographique}

Quatre composantes : régularité hydrique (35 points), vigueur végétative (25), santé du sol
(20), intégrité (20). Une chaîne invalide, un scellé rompu ou un écart satellite supérieur
à 15 points annulent la composante intégrité et suspendent le décaissement. Le dossier
complet s'exporte en PDF numéroté, accompagné d'une \textbf{archive ZIP contenant un
vérificateur autonome} qui recalcule les empreintes en JavaScript pur, sans connexion.

\subsubsection{Trois espaces, un seul jeu de données}

L'exploitant pilote sa parcelle et consent au partage ; la banque consulte le dossier
seulement après ce consentement ; l'agent de conformité certifie ou rejette les dossiers
avec note obligatoire, horodatage et identité du relecteur. Aucune vue ne recalcule les
règles : le moteur de score est unique et partagé.

% ══════════════════════════════════════════════════════════
\newpage
\section{Architecture technique}
% ══════════════════════════════════════════════════════════

\subsection{Chaîne de la preuve}

\begin{center}
\begin{tikzpicture}[
  b/.style={
    rectangle, rounded corners=4pt, minimum width=2.9cm, minimum height=1.7cm,
    draw=primary, fill=primary!8, thick, align=center,
    font=\footnotesize\bfseries\color{primary}
  },
  h/.style={
    rectangle, rounded corners=4pt, minimum width=2.9cm, minimum height=1.7cm,
    draw=techblue, fill=techblue!8, thick, align=center,
    font=\footnotesize\bfseries\color{techblue}
  },
  arr/.style={-{Stealth[length=5pt]}, thick, muted}
]
  \node[b] (r1) {Relevé $n-1$\\ {\tiny mesures + scellé}};
  \node[h, right=0.7cm of r1] (h1) {SHA-256\\ {\tiny hash $n-1$}};
  \node[b, right=0.7cm of h1] (r2) {Relevé $n$\\ {\tiny mesures + scellé}};
  \node[h, right=0.7cm of r2] (h2) {SHA-256\\ {\tiny hash $n$}};
  \node[b, below=1.0cm of h2] (score) {Score Vert\\ {\tiny 4 composantes}};
  \node[b, left=0.7cm of score] (pdf) {Dossier PDF\\ {\tiny n° AT-…}};
  \node[b, left=0.7cm of pdf] (bank) {Portail prêteur\\ {\tiny consentement}};

  \draw[arr] (r1) -- (h1);
  \draw[arr] (h1) -- (r2);
  \draw[arr] (r2) -- (h2);
  \draw[arr] (h2) -- (score);
  \draw[arr] (score) -- (pdf);
  \draw[arr] (pdf) -- (bank);
\end{tikzpicture}
\end{center}

\subsection{Stack du démonstrateur}

\begin{tabularx}{\linewidth}{>{\bfseries\color{primary}}l l X}
  \toprule
  \rowcolor{lightbg}
  Composant & Technologie & Justification \\
  \midrule
  Interface & React 18 + TypeScript 5.6 & Typage strict, composants partagés \\
  \addlinespace
  Build & Vite 5.4 & Démarrage rapide, sortie statique \\
  \addlinespace
  Style & Tailwind CSS 3.4 & Palette, responsive 375\,px, contrastes \\
  \addlinespace
  Cartographie & Leaflet + react-leaflet & Fond satellite, contours, sondes \\
  \addlinespace
  Géométrie & Turf.js & Surface, centroïde, découpage en zones \\
  \addlinespace
  Courbes & Recharts & Télémétrie sur 90 jours \\
  \addlinespace
  Preuve & Web Crypto (SHA-256) & Chaînage sans dépendance externe \\
  \addlinespace
  Stockage & localStorage + IndexedDB & Structure et pièces, sans serveur \\
  \addlinespace
  Export & jsPDF + autotable, ZIP maison & Dossier de crédit et archive vérifiable \\
  \bottomrule
\end{tabularx}

\subsection{Ce que le prototype ne fait pas}

\begin{itemize}
  \item Aucune connexion à un boîtier physique : les 90 jours de relevés sont générés par
        des fonctions déterministes, et les couches Sentinel sont simulées.
  \item Aucun backend : les comptes et les dossiers restent sur l'appareil ; l'archive ZIP
        est la seule copie portable.
  \item La signature prouve l'absence de réécriture locale, pas l'origine physique de la
        mesure : c'est le rôle du scellé matériel et de la passerelle.
\end{itemize}

% ══════════════════════════════════════════════════════════
\section{Éléments justifiant le caractère innovant}
% ══════════════════════════════════════════════════════════

\begin{enumerate}
  \item \textbf{La preuve est le produit, pas un rapport.} Les plateformes agricoles
        existantes affichent des courbes ; AgriTwin produit un objet vérifiable : chaîne
        de hachage, score décomposé, PDF numéroté, vérificateur hors ligne. La banque ne
        fait pas confiance à l'écran : elle recalcule.

  \item \textbf{La fraude est traitée comme un cas nominal.} Le rapprochement
        sonde / satellite (écart supérieur à 15 points), le scellé et l'accéléromètre sont
        des entrées du score, pas des alertes décoratives. Une fraude ne dégrade pas la
        note : elle annule la composante intégrité et gèle le décaissement.

  \item \textbf{Le score est auditable par conception.} Poids, seuils, formules et limites
        sont publiés dans l'application et importés du même moteur que le calcul : il n'y a
        pas deux versions de la vérité, une pour l'écran et une pour le règlement.

  \item \textbf{Le coût marginal de la donnée est faible.} Le boîtier est un investissement
        unique (≈ 300\,000 FCFA pour une parcelle de 1,5\,ha), le satellite est gratuit,
        et l'application ne dépend d'aucun service payant. Le modèle économique ne repose
        pas sur un abonnement de données coûteux.

  \item \textbf{Le consentement est un mécanisme, pas une case.} Sans consentement explicite
        de l'exploitant, le portail prêteur est verrouillé ; le partage est révocable et
        tracé, conformément à la loi 2008-12 sur les données personnelles.

  \item \textbf{L'inclusion linguistique est intégrée.} Alertes, questions rapides et
        restitutions existent en français et en wolof, avec un repère sonore, pour un
        utilisateur qui lit difficilement le français technique.
\end{enumerate}

\begin{mdframed}[style=neutral]
\textbf{Positionnement face aux approches existantes.} Les stations météo et les
plateformes d'irrigation automatisée mesurent et pilotent ; elles ne produisent pas de
preuve bancaire et coûtent de 650\,000 FCFA à plus de 2\,M FCFA installées. Les
applications de conseil agricole conseillent sans historiser de façon opposable. AgriTwin
ne vend ni la mesure ni l'automatisation : il vend \textbf{la convertibilité d'une parcelle
en dossier de crédit}.
\end{mdframed}

% ══════════════════════════════════════════════════════════
\newpage
\section{Niveau de développement du projet}
% ══════════════════════════════════════════════════════════

\subsection{Un prototype fonctionnel, pas une maquette}

Le démonstrateur est une application monopage complète, déployable en statique, qui
implémente l'ensemble du parcours :

\begin{itemize}
  \item authentification locale avec trois rôles et trois comptes de démonstration ;
  \item déclaration de parcelle en quatre étapes, tracé du contour, contrôle binaire des
        pièces et stockage IndexedDB ;
  \item carte satellite, trois calques (zones, NDVI, radar), sondes cliquables et curseur
        temporel sur 90 jours ;
  \item télémétrie, journal agronomique, bilan azoté, assistant vocal et alertes bilingues ;
  \item simulateur de rendement à trois leviers, avec marge nette en FCFA ;
  \item moteur de Score Vert à quatre composantes, verdict de crédit et assurance indicielle ;
  \item chaîne SHA-256 vérifiable, journal de blocs, PDF de crédit, archive ZIP et
        vérificateur autonome ;
  \item portail financeur verrouillé par le consentement, back-office de conformité ;
  \item simulateur d'injection (relevé normal, chaleur, fraude, arrachement) et démo guidée
        en neuf étapes ;
  \item mise en page vérifiée de 320 à 375\,px, navigation clavier, libellés ARIA.
\end{itemize}

\subsection{État de maturité}

\begin{center}
\begin{tabularx}{\linewidth}{>{\bfseries\color{primary}}l X}
  \toprule
  \rowcolor{lightbg}
  Dimension & État au dépôt \\
  \midrule
  Logiciel & Fonctionnel, testé manuellement, build de production reproductible \\
  \addlinespace
  Données & Simulées de façon déterministe (90 jours, 3 zones, 3 parcelles de démonstration) \\
  \addlinespace
  Matériel & Spécifié et chiffré (4 configurations de kits), non encore assemblé \\
  \addlinespace
  Terrain & Aucun pilote réel à ce stade : c'est la prochaine étape \\
  \addlinespace
  Niveau TRL & 4 — validation en environnement de laboratoire / démonstrateur \\
  \bottomrule
\end{tabularx}
\end{center}

\begin{mdframed}[style=accent]
\textbf{Honnêteté requise devant le jury.} Le prototype ne prétend ni mesurer un sol réel,
ni accorder un crédit. Il démontre que la chaîne « mesure → preuve → score → décision »
fonctionne de bout en bout, et que chaque maillon est vérifiable indépendamment.
\end{mdframed}

% ══════════════════════════════════════════════════════════
\section{Perspectives d'application et de valorisation}
% ══════════════════════════════════════════════════════════

\subsection{Segments cibles}

\begin{tabularx}{\linewidth}{>{\bfseries\color{primary}}l X l}
  \toprule
  \rowcolor{lightbg}
  Segment & Besoin principal & Offre \\
  \midrule
  Petit producteur (0,5–2\,ha) & Prouver sa parcelle pour un crédit d'intrants & Kit A + plateforme \\
  \addlinespace
  Exploitation moyenne (2–10\,ha) & Suivre trois zones, justifier les rendements & Kit B + tableau de bord \\
  \addlinespace
  Coopérative (20–100\,ha) & Dossier collectif, passerelle mutualisée & Kit C + portail coopérative \\
  \addlinespace
  IMF / banque & Instruire à distance, réduire le risque & Licence portail + API \\
  \addlinespace
  Assureur indiciel & Déclencher une indemnisation sur des données vérifiables & Licence + connecteur \\
  \addlinespace
  Recherche / agro-industrie & Métrologie de référence, essais & Kit D + export de données \\
  \bottomrule
\end{tabularx}

\subsection{Modèle de valorisation}

\begin{tabularx}{\linewidth}{>{\bfseries}l X >{\color{mygreen}\small\itshape}r}
  \toprule
  \rowcolor{lightbg}
  Axe & Description & Modèle \\
  \midrule
  \multicolumn{3}{l}{\textbf{\textcolor{primary}{Matériel}}} \\
  \addlinespace
  Kit A / B / C & Boîtier scellé, sondes, énergie, pose & 300\,000 à 2\,100\,000 FCFA \\
  Leasing (HaaS) & Apport 30\,\%, puis 24 à 36 mois & 15\,000 à 25\,000 FCFA/mois \\
  \addlinespace\midrule\addlinespace
  \multicolumn{3}{l}{\textbf{\textcolor{primary}{Plateforme}}} \\
  \addlinespace
  Abonnement exploitation & Score, alertes, PDF, archive & 5\,000 à 15\,000 FCFA/mois \\
  Licence institution & Portail prêteur, API, dossiers illimités & Contrat annuel \\
  Intégration coopérative & Formation, support, hébergement & Forfait par coopérative \\
  \bottomrule
\end{tabularx}

\subsection{Feuille de route}

\begin{tabularx}{\linewidth}{>{\bfseries\color{primary}}l l X}
  \toprule
  \rowcolor{lightbg}
  Phase & Horizon & Objectifs \\
  \midrule
  Phase 1 & 0–3 mois & Assembler 3 kits A, valider la sonde 4-en-1, première parcelle pilote à Bambey \\
  \addlinespace
  Phase 2 & 3–9 mois & Backend et synchronisation, 10 parcelles, dossier réel présenté à une IMF \\
  \addlinespace
  Phase 3 & 9–18 mois & Passerelle coopérative, 100 parcelles, connecteur assureur indiciel \\
  \addlinespace
  Phase 4 & 18 mois+ & Extension Thiès / Mbour / Saint-Louis, API opérateurs \\
  \bottomrule
\end{tabularx}

% ══════════════════════════════════════════════════════════
\newpage
\section{Impacts économiques, sociaux et environnementaux attendus}
% ══════════════════════════════════════════════════════════

\begin{tabularx}{\linewidth}{>{\bfseries\color{primary}}p{3.1cm} X X}
  \toprule
  \rowcolor{lightbg}
  Dimension & Mécanisme & Indicateur de suivi \\
  \midrule
  \textbf{Économique} &
  Dossier instruisible à distance : moins de déplacements d'experts, décision plus rapide,
  crédit d'intrants mieux calibré. &
  Délai d'instruction (jours), coût de vérification par dossier, taux de dossiers complets \\
  \addlinespace
  \textbf{Économique} &
  Le score décomposé permet de négocier une bonification ou une assurance adaptée au risque réel. &
  Taux effectif obtenu, part des dossiers à score A \\
  \addlinespace
  \textbf{Social} &
  Inclusion des producteurs sans historique bancaire, y compris les femmes des coopératives,
  grâce à un dossier construit sur des faits mesurés. &
  Nombre de primo-emprunteurs, part de femmes, rétention à 12 mois \\
  \addlinespace
  \textbf{Social} &
  Alertes et restitutions en wolof, navigation mobile : l'outil reste utilisable par un
  producteur peu lettré en français technique. &
  Taux d'usage effectif des alertes, tests d'utilisabilité \\
  \addlinespace
  \textbf{Environnemental} &
  Irrigation raisonnée sur la fenêtre 20–35\,\% d'humidité et azote ajusté au besoin de la
  culture : moins de gaspillage d'eau et d'engrais. &
  Volume d'eau par hectare, kg N/ha, nombre d'alertes de stress évitées \\
  \addlinespace
  \textbf{Environnemental} &
  Preuve partagée numériquement : moins de tournées de contrôle et de dossiers papier. &
  Déplacements évités, volume de pièces dématérialisées \\
  \addlinespace
  \textbf{Institutionnel} &
  Journal opposable et consentement tracé : la décision devient défendable en cas de
  contestation ou d'audit. &
  Nombre de dossiers audités sans rupture de chaîne \\
  \bottomrule
\end{tabularx}

\vspace{0.5cm}

\begin{mdframed}[style=primary]
\textbf{Prudence méthodologique.} Aucun de ces impacts n'est mesuré à ce jour : le
prototype n'a pas été déployé en champ. Les valeurs seront établies lors du pilote par
comparaison avant / après sur un échantillon témoin, et publiées telles quelles, y compris
si elles sont inférieures aux attentes.
\end{mdframed}

% ══════════════════════════════════════════════════════════
\section{Risques, limites et éthique}
% ══════════════════════════════════════════════════════════

\begin{tabularx}{\linewidth}{>{\bfseries\color{primary}}p{3.4cm} X X}
  \toprule
  \rowcolor{lightbg}
  Risque & Description & Réduction \\
  \midrule
  Métrologie & Sondes bas coût : pH et conductivité indicatifs, dérive possible. & Comparaison contre une sonde de référence avant toute promesse de précision ; mention explicite dans l'interface. \\
  \addlinespace
  Connectivité & Couverture LoRa ou data insuffisante en zone rurale. & Relevés en lot, cache local, passerelle mutualisée, mode hors ligne assumé. \\
  \addlinespace
  Sécurité & Le stockage navigateur est modifiable par un utilisateur averti. & Le prototype est un démonstrateur ; la version pilote déplace la preuve côté serveur et conserve le vérificateur indépendant. \\
  \addlinespace
  Données personnelles & Identité, pièces, localisation de la parcelle. & Consentement explicite, partage révocable, référence à la loi 2008-12 (CDP), minimisation des données. \\
  \addlinespace
  Adoption & Coût du kit et fracture numérique. & Leasing, subvention coopérative, interface en wolof, formation de deux heures à la pose. \\
  \addlinespace
  Responsabilité & Un score ne doit pas se substituer à la décision de crédit. & Le score est un \emph{avis documenté} : la banque reste décideur, le motif de refus est explicite. \\
  \bottomrule
\end{tabularx}

% ══════════════════════════════════════════════════════════
\section{Conclusion}
% ══════════════════════════════════════════════════════════

AgriTwin part d'un constat simple : une parcelle bien conduite ne vaut rien pour une banque
tant que rien ne le prouve. Plutôt que d'ajouter une énième application de conseil, le
projet construit la \textbf{preuve} elle-même — mesure scellée, empreinte chaînée, score
décomposé, dossier vérifiable hors ligne — et la rend lisible par les trois parties
concernées : le producteur, le financeur, le contrôleur.

Le prototype démontre cette chaîne de bout en bout, avec ses limites assumées : données
simulées, matériel spécifié mais non assemblé, aucun pilote terrain. La prochaine étape est
donc claire et vérifiable : trois kits, dix parcelles, un dossier réel présenté à une
institution de microfinance, et des indicateurs publiés sans embellissement.

\vspace{0.4cm}

Face aux enjeux du financement de l'agriculture familiale sénégalaise, AgriTwin incarne ce
que l'ingénierie peut apporter de plus utile : \textbf{une preuve simple à produire,
difficile à falsifier, et immédiatement utilisable par ceux qui décident}.

\vspace{1.2cm}

\begin{center}
  \textcolor{rule}{\rule{4cm}{0.4pt}}
  \quad\textcolor{accent}{\large$\bullet$}\quad
  \textcolor{rule}{\rule{4cm}{0.4pt}}

  \vspace{0.5cm}
  {\Large\textbf{\textcolor{primary}{AGRITWIN}}}\\[5pt]
  \textit{\textcolor{muted}{Une parcelle, une preuve, un score.}}\\[8pt]
  \small\textcolor{muted}{Prix de l'Innovation 2026 \quad—\quad Polytech Innovation Day \quad—\quad EPT Thiès}
\end{center}

% ══════════════════════════════════════════════════════════
\newpage
\appendix
\section{Méthodologie du Score Vert}
% ══════════════════════════════════════════════════════════

\begin{center}
\begin{tabularx}{\linewidth}{>{\bfseries\color{primary}}l c X}
  \toprule
  \rowcolor{lightbg}
  Composante & Poids & Calcul \\
  \midrule
  Régularité hydrique & 35 & Part des jours où l'humidité à 15\,cm reste dans la fenêtre 20–35\,\% \\
  \addlinespace
  Vigueur végétative & 25 & NDVI multiplié par 32, plafonné à 25 \\
  \addlinespace
  Santé du sol & 20 & 20 points si EC $\leq$ 1,4\,dS/m et pH entre 6,0 et 7,5 ; 14 si EC $\leq$ 2,0 ; sinon 7 \\
  \addlinespace
  Intégrité & 20 & 20 points si chaîne valide, scellé intact et écart satellite $\leq$ 15 ; sinon 0 \\
  \bottomrule
\end{tabularx}
\end{center}

Seuils de décision : \textbf{75 et plus} profil A (éligible), \textbf{50 à 74} profil B
(suivi coopératif), \textbf{moins de 50} profil C. Une fraude, une chaîne invalide ou un
dossier gelé suspend le décaissement, quel que soit le score affiché.

\section{Composition et coût des kits}
% ══════════════════════════════════════════════════════════

\begin{center}
\begin{tabularx}{\linewidth}{>{\bfseries\color{primary}}l c X r}
  \toprule
  \rowcolor{lightbg}
  Kit & Surface & Composition principale & Coût total \\
  \midrule
  A — Essentiel & 0,5–2\,ha & 2 sondes 4-en-1, nœud LoRa, boîtier IP67, 20\,Wc, pose & $\approx$ 300\,000 FCFA \\
  \addlinespace
  B — Trois zones & 2–10\,ha & 3 sondes, boîtier 255\,mm, 40\,Wc, bus RS485 & $\approx$ 450\,000 FCFA \\
  \addlinespace
  C — Coopérative & 20–100\,ha & 3 kits B, passerelle extérieure, station météo & $\approx$ 2\,100\,000 FCFA \\
  \addlinespace
  D — Recherche & 1–5\,ha & Sondes de référence, passerelle, étalonnage & $\approx$ 1\,300\,000 FCFA \\
  \bottomrule
\end{tabularx}
\end{center}

Prix indicatifs constatés en novembre 2026, matériel rendu Dakar, pose comprise ; à
remplacer par des devis fournisseurs avant tout engagement. Coût de possession sur trois
ans pour le kit A : environ 620\,000 FCFA, abonnement et maintenance inclus.

\section{Protocole de vérification indépendante}
% ══════════════════════════════════════════════════════════

L'archive remise à la banque contient le PDF du dossier, un fichier \texttt{preuve.json}
(format \texttt{agritwin-dossier/1}) et un vérificateur HTML autonome. Ce vérificateur
réalise six contrôles, sans réseau et sans bibliothèque : format du dossier, chaîne des
blocs, empreinte globale, recalcul du Score Vert et de ses quatre composantes, empreinte du
PDF déposé, cohérence de l'identité. Toute altération est localisée : le bloc de rupture et
l'écart entre score déclaré et score recalculé sont affichés.

\section{Check-list de dépôt}
% ══════════════════════════════════════════════════════════

\begin{itemize}
  \item Renseigner la page de garde : noms, formations, encadrant, contact.
  \item Insérer le logo à la place du cadre vide, puis vérifier la mise en page.
  \item Remplacer l'URL du prototype par l'adresse publique réelle.
  \item Joindre les supports autorisés : photos de parcelle et de boîtier, schéma
        d'architecture, capture mobile, vidéo de démonstration.
  \item Relire les chiffres de coût et les dater (« prix constatés en novembre 2026 »).
  \item Vérifier que chaque affirmation sur les impacts est présentée comme une hypothèse
        à mesurer, jamais comme un résultat acquis.
\end{itemize}

\end{document}
